\section{課題解決のための提案}

前章では、eスポーツの市場規模やファン人口が増加している一方で、
地域やゲームコミュニティなどが開催する大会では、競技者情報や試合結果の管理、
情報共有に手作業が発生しやすいことなどを課題としてあげました。また、eスポーツ大会では、
競技者だけではなく観戦者に対しても、大会の進行状況や試合結果などの情報を提供することが重要です。
特に、注目選手が参加する大会では、その選手の試合状況や今後の観戦予定を確認する観戦者も存在します。

そこで、本システムでは大会運営に必要な情報を一元管理し、主催者・競技者・観戦者がそれぞれ必要な情報を確認できる
「小規模ゲーム大会運営システム(Play Line)」を提案します。本システムでは、競技者の参加登録から
試合結果の管理までを効率化するとともに、大会情報とライブ配信へのアクセスを関連付けることで、
競技者が自分の試合情報を確認しやすくし、観戦者が大会全体や注目選手の試合を追いかけやすい環境の実現を目指します。

\subsection{小規模ゲーム大会運営システム(Play Line)}

本システムでは、大会の参加登録から大会終了までに必要となる競技者情報、対戦情報、試合結果などを
データベースで一元管理します。これにより、DiscordやExcel、Googleスプレッドシートなどに分散していた
大会日程、参加者情報、対戦表、試合結果などを一元管理し、主催者の運営負担を軽減するとともに、
参加者や観戦者が必要な情報を容易に確認できるようにします。

まず、主催者が本システムに大会情報や参加条件などを登録し、競技者の募集を行います。
競技者はWebブラウザから本システムのWebサイトへアクセスし、参加したい大会の登録を行います。
登録された情報はデータベースに保存され、主催者は登録された競技者の情報を一覧で確認できるようにします。

大会開催時には、登録された競技者情報を利用してトーナメント表を作成します。
競技者はWebサイトから自分の対戦相手やトーナメント上の位置を確認できます。
また、主催者は試合終了後に試合結果を登録し、その結果に基づいて大会の進行状況を更新します。
これにより、競技者は次の対戦相手や次の試合の確認できるようになります。
主催者が登録・更新した大会情報や試合結果はデータベースに保存され、観戦者にも提供されます。
観戦者は大会ページから、トーナメント表、試合結果、現在行われている試合、今後行われる試合などを確認できます。

本システムにおける主催者による情報の登録・更新から、競技者や観戦者への情報提供までの流れを以下に示します。

\begin{figure}[H]
  \centering
  \begin{tikzpicture}[
    node distance=0.9cm,
    box/.style={
      rectangle, rounded corners=4pt,
      draw=blue!60!black, fill=blue!8,
      minimum width=7cm, minimum height=0.9cm,
      align=center, thick
    },
    arrow/.style={-{Stealth[length=2.5mm]}, thick, draw=blue!60!black}
  ]
    \node[box] (n1) {競技者受付};
    \node[box, below=of n1] (n2) {競技者情報管理};
    \node[box, below=of n2] (n3) {トーナメント作成};
    \node[box, below=of n3] (n4) {対戦相手確認};
    \node[box, below=of n4] (n5) {試合結果登録};
    \node[box, below=of n5] (n6) {大会進行状況更新};
    \node[box, below=of n6] (n7) {競技者・観戦者への情報提供};

    \foreach \i/\j in {1/2,2/3,3/4,4/5,5/6,6/7} {
      \draw[arrow] (n\i) -- (n\j);
    }
  \end{tikzpicture}
  \caption{大会運営の業務フロー}
  \label{fig:flow}
\end{figure}

このように、大会運営に必要な情報を一つのシステムで管理することで、主催者による
情報の整理や更新の負担を軽減します。また、主催者が登録・更新した大会情報を
競技者と観戦者が同じWebサイトから確認できるようにすることで、個別に情報を伝達する手間を減らし、
対戦相手や試合結果などの情報を効率的に共有できるようにします。
これにより、利用者が必要な情報を容易に取得できる仕組みの実現を目指します。

\subsection{注目選手の参加を活かした観戦者向け情報提供}

1章で述べたように、注目選手の参加は大会への関心を高める要因の一つとなる可能性がある。
そこで本システムでは、大会運営の効率化に加え、観戦者が関心のある選手の試合を継続して
追いやすくするための情報提供を行います。例えば、大会ページから

\begin{itemize}
  \item 注目選手（運営による注目）の現在の対戦状況
  \item 注目選手の過去の全ての試合結果
  \item 次の対戦相手
  \item トーナメント上の現在位置
  \item 大会全体の進行状況
  \item 各試合の結果
\end{itemize}
などを確認できるようにします。

これにより、特定の選手を応援している観戦者が、その選手の現在の状況や次の試合を把握しやすくなります。
また、地域大会において実績のある選手や地域で人気のある選手が参加した場合にも、
その選手の対戦状況や試合結果をシステム上で確認できるようにすることで、
会場やオンラインで大会を観戦するきっかけを作り、大会への関心を高めることを目指します。

eスポーツの観戦動機に関する研究では、1,052人を対象とした調査により、好きな選手への
憧れや応援等の動機が、ライブ配信の視聴行動に関係することが報告されています。\cite{scientificreports2024}
\footnote{Yuxin Tang, Xiaowen Zhang, Shengfeng Zan, “Exploring e-sports fans’ motivation for watching live streams based on self-determination theory,” Scientific Reports, Vol. 14, Article 13858, 2024. 

自己決定理論に基づくeスポーツファンのライブ配信視聴動機を探る |科学報告書　（閲覧日：2026年10月15日） }

このことから、観戦者が関心のある選手の試合情報を確認できることは、試合への関心を高め、
観戦を楽しむための一助になると考えられます。
そこで、本システムでは観戦者が選手の情報や現在の対戦状況、過去の試合結果、次の対戦相手、
トーナメント表における現在の位置などを大会ページから確認できるようにします。
これにより、観戦者が注目選手を探す手間を減らし、注目選手やお気に入りの選手の試合状況を把握しやすくすることを目指します。

\subsection{競技者と観戦者が大会状況を確認できる仕組み}

2.1では大会運営に必要な情報を一元管理する仕組みについて、2.2では選手に関する情報を観戦者に提供する仕組みについて述べました。
これらを実現する為、本システムでは、利用者の目的に応じて必要な大会情報を確認できる画面を提供します。

競技者には、自分の試合に関する情報を中心に表示します。競技者はWebサイトから、
参加登録状況、自分の対戦相手、試合結果、トーナメント表における現在の位置、次の試合などを確認できるようにします。
これにより、競技者が自分の試合に必要な情報を簡単に取得できるようにします。

一方、観戦者には大会全体の状況や注目選手の情報を確認できる画面を提供します。
具体的には、現在行われている試合、各試合の結果、トーナメント表、今後行われる試合、注目選手の対戦状況や
過去の試合結果、大会全体の進行状況などを表示します。
更に、本システムでは大会のライブ配信URLを大会情報に登録し、観戦者が大会ページから配信に直接アクセスできる仕組みを設けます。
主催者は、大会の作成時や編集時にライブ配信URLを登録し、観戦者は大会ページに表示されたリンクから
配信ページに移動できるようにします。これにより、観戦者が配信URLを別のWebサイトやコミュニケーションツール等で探す手間を軽減します。

以上の仕組みにより、本システムでは主催者の運営負担を軽減するとともに、競技者が必要な試合情報を確認しやすくし、
観戦者が大会情報とライブ配信を関連付けて利用できる環境の実現を目指します。

\subsection{βテストによる効果の検証}

本システムの有用性については、実際のゲーム大会競技者、観戦者、ゲームコミュニティなど
開発者以外の利用者を対象としたβテストを実施して検証します。
βテストでは、主催者・競技者・観戦者の3者に実際にシステムを利用してもらい、
それぞれの立場からシステムの効果を評価します。

主催者には、競技者受付、競技者情報の確認、トーナメント作成、試合結果の登録、大会情報の管理などを行ってもらい、
以下の作業時間を計測します。

\begin{itemize}
  \item 競技者受付にかかる時間
  \item 競技者情報を整理する時間
  \item トーナメントを作成する時間
  \item 試合結果を登録する時間大会全体の情報管理にかかる時間
\end{itemize}

また、競技者は以下の操作を行ってもらいます。

\begin{itemize}
  \item 大会の参加登録から大会終了までに必要となる競技者情報自分の対戦相手の確認
  \item 試合結果の確認
  \item 次の試合の確認
\end{itemize}

同様に、観戦者には以下の操作を行ってもらいます。

\begin{itemize}
  \item トーナメント表の確認
  \item 現在の試合状況の確認
  \item 試合結果の確認
  \item 大会全体の進行状況の確認
  \item 注目選手の対戦状況の確認(運営による注目選手の設定)
\end{itemize}

更に、主催者・競技者・観戦者それぞれにアンケートを実施し、
\begin{itemize}
  \item 操作のわかりやすさ
  \item 大会情報の確認のしやすさ
  \item 登録作業のしやすさ
  \item 対戦相手の確認のしやすさ
  \item 大会結果の確認のしやすさ
  \item 大会進行状況のわかりやさ宇佐
  \item 注目選手の試合を追いやすいか
  \item 大会を観戦する上で役立つか
  \item 大会運営の負担が軽減されたと感じるか
  \item 大会ページからライブ配信を見つけやすかったか
  \item 大会情報を確認しながらライブ配信を利用しやすかったか
\end{itemize}
などを調査します。

主催者については、システム導入前後の作業時間を比較し、

\begin{equation}
  \text{削減率} = \frac{\text{導入前の時間} - \text{導入後の時間}}{\text{導入前の時間}} \times 100 \ [\%]
  \label{eq:reduction_rate}
\end{equation}

として評価します。

また、競技者については、自分の対戦相手や試合結果を確認するまでにかかる時間などを比較し、
試合情報のしやすさを評価します。観戦者については、大会の進行状況や注目選手の試合を追いやすくなったかを
アンケートや情報確認にかかる時間によって評価します。

このように、βテストでは単に「システムが動作するか」を確認するのではなく、
主催者・競技者・観戦者のそれぞれにどのような効果があるのかを検証します。

以上の提案により、本システムでは、大会運営に必要な情報を一元的に管理することで
主催者の作業負担を軽減し、参加者には自分の試合に必要な情報を提供し、
観戦者には大会全体や注目選手の試合を追いやすい環境を提供します。

これにより、小規模なゲーム大会を単に「ゲームをプレイする場」とするのではなく、
主催者が運営しやすく、参加者が競技しやすく、観戦者も大会を楽しめる場として発展させることを目指します。
